\documentclass[11pt,a4paper]{article}
\usepackage{auditstyle}

\newcommand{\eps}{\varepsilon}
\newcommand{\W}{\mathcal W}
\newcommand{\dd}{\,d}
\newtheorem{corollary}[theorem]{Corollary}

\title{Part II static projective compactness\\[3pt]
\large A fixed-coarse interacting scaling state along a cofinal subnet}
\author{Lamport-form proof}
\date{9 August 2026}

\begin{document}
\maketitle

\begin{resultbox}
\textbf{Unconditional static conclusion.}  Once one state has been chosen on every
finite-resolution algebra, weak-star compactness produces a cofinal subnet whose
pullbacks converge simultaneously on every fixed coarse algebra.  The limiting family
is projectively consistent and therefore defines a state on the \(C^*\)-inductive limit.

The theorem proves neither convergence of the full sequence nor identification with the
continuum \(P(\phi)_2\) vacuum.  It gives no rate, no local-state-norm convergence, no
normality, and no limiting dynamics.
\end{resultbox}

\tableofcontents

\section{Statement}

\begin{definition}[directed algebraic system]
For \(M\le N\), let \(\A_M,\A_N\) be unital \(C^*\)-algebras and let
\[
 \alpha_M^N:\A_M\longrightarrow\A_N                              \tag{1.1}
\]
be injective unital \( * \)-homomorphisms such that
\[
 \alpha_N^K\circ\alpha_M^N=\alpha_M^K,\qquad M\le N\le K.        \tag{1.2}
\]
Let \((\A_\infty,\alpha_M^\infty)\) be their \(C^*\)-inductive limit, so
\[
 \alpha_N^\infty\circ\alpha_M^N=\alpha_M^\infty,\qquad
 \overline{\bigcup_M\alpha_M^\infty(\A_M)}^{\|\cdot\|}
 =\A_\infty.                                                    \tag{1.3}
\]
\end{definition}

\begin{theorem}[cofinal-subnet projective limit]\label{thm:compactness}
For every \(N\), let \(\omega_N\) be a state on \(\A_N\).  Then there exist a
directed set \(I\), a cofinal subnet map \(i\mapsto N(i)\in\mathbb N_0\), states
\(\rho_M\in S(\A_M)\), and a state \(\Omega_\infty\in S(\A_\infty)\) such that
\begin{align}
 \lim_{i\in I}\omega_{N(i)}(\alpha_M^{N(i)}(A))
 &=\rho_M(A)
 =\Omega_\infty(\alpha_M^\infty(A)),                            \tag{1.4}\\
 \rho_M&=\rho_{M'}\circ\alpha_M^{M'},\qquad M\le M',             \tag{1.5}
\end{align}
for every fixed \(M\) and every \(A\in\A_M\).
\end{theorem}

\section{Proof of Theorem \ref{thm:compactness}}

\LS{1}{1}{Prove weak-star compactness at one scale.}

The closed unit ball of \(\A_M^*\) is compact in
\(\sigma(\A_M^*,\A_M)\) by Banach--Alaoglu.  Within that ball,
\[
 S(\A_M)=
 \{\varphi:\varphi(1)=1,\ \varphi(A^*A)\ge0\text{ for every }A\in\A_M\}. \tag{2.1}
\]
For fixed \(A\), both maps
\[
 \varphi\longmapsto\varphi(1),\qquad
 \varphi\longmapsto\varphi(A^*A)                               \tag{2.2}
\]
are weak-star continuous.  Hence every condition in (2.1) is weak-star closed.
Therefore \(S(\A_M)\) is weak-star compact.

\LS{1}{2}{Put all coarse restrictions into one compact space.}

Define
\[
 K:=\prod_{M=0}^{\infty}S(\A_M),                               \tag{2.3}
\]
with the product of the weak-star topologies.  Tychonoff's theorem makes \(K\)
compact.  For every \(N\), define \(x_N=(x_N^M)_{M\ge0}\in K\) by
\[
 x_N^M:=
 \begin{cases}
  \omega_N\circ\alpha_M^N,&M\le N,\\
  \omega_M,&M>N.
 \end{cases}                                                    \tag{2.4}
\]
The second case only fills coordinates that are eventually irrelevant.

\LS{1}{3}{Extract a genuinely cofinal subnet.}

Every net in a compact space has a convergent subnet.  Hence there are a directed set
\(I\), a subnet map \(N:I\to\mathbb N_0\), and
\((\rho_M)_M\in K\) such that
\[
 x_{N(i)}\longrightarrow(\rho_M)_M\quad\hbox{in }K.             \tag{2.5}
\]
By the definition of a subnet, the map \(i\mapsto N(i)\) is cofinal: for every
\(M'\), there is \(i_0\) such that \(N(i)\ge M'\) whenever \(i\ge i_0\).
Product convergence in (2.5) means, explicitly,
\[
 \lim_i x_{N(i)}^M(A)=\rho_M(A),\qquad A\in\A_M.                \tag{2.6}
\]

\LS{1}{4}{Prove projective consistency.}

Fix \(M\le M'\) and \(A\in\A_M\).  Cofinality implies that \(N(i)\ge M'\)
eventually.  For those \(i\), (1.2) gives
\begin{align}
 x_{N(i)}^{M'}(\alpha_M^{M'}(A))
 &=\omega_{N(i)}\bigl(\alpha_{M'}^{N(i)}
      (\alpha_M^{M'}(A))\bigr)\\
 &=\omega_{N(i)}(\alpha_M^{N(i)}(A))
 =x_{N(i)}^M(A).                                                \tag{2.7}
\end{align}
Take the limit in (2.7) and use (2.6):
\[
 \rho_{M'}(\alpha_M^{M'}(A))=\rho_M(A).                         \tag{2.8}
\]
This is exactly (1.5).

\LS{1}{5}{Define the functional on the algebraic inductive union.}

Set
\[
 \A_{\rm alg}:=\bigcup_M\alpha_M^\infty(\A_M)                  \tag{2.9}
\]
and prescribe
\[
 \Omega_0(\alpha_M^\infty(A)):=\rho_M(A).                      \tag{2.10}
\]
To prove that (2.10) is well-defined, suppose
\[
 \alpha_M^\infty(A)=\alpha_{M'}^\infty(B).                     \tag{2.11}
\]
Choose \(K\ge M,M'\).  By (1.3) and injectivity of
\(\alpha_K^\infty\),
\[
 \alpha_M^K(A)=\alpha_{M'}^K(B).                               \tag{2.12}
\]
Applying \(\rho_K\) and then (2.8) gives
\[
 \rho_M(A)=\rho_K(\alpha_M^K(A))
 =\rho_K(\alpha_{M'}^K(B))=\rho_{M'}(B).                       \tag{2.13}
\]
Thus (2.10) is independent of the representative.

\LS{1}{6}{Prove positivity, normalization, and boundedness.}

Every finite collection of elements of \(\A_{\rm alg}\) lies in the image of one
\(\A_K\).  Consequently linearity of \(\Omega_0\) follows from linearity of
\(\rho_K\).  If \(C=\alpha_M^\infty(A)\), then
\[
 \Omega_0(C^*C)=\rho_M(A^*A)\ge0,\qquad \Omega_0(1)=1.          \tag{2.14}
\]
For completeness, positivity implies the Cauchy--Schwarz inequality
\[
 |\rho_M(A)|^2\le\rho_M(A^*A)\rho_M(1)
 \le\|A\|^2.                                                    \tag{2.15}
\]
Because every injective \( * \)-homomorphism of \(C^*\)-algebras is isometric,
\[
 |\Omega_0(\alpha_M^\infty(A))|
 \le\|A\|=\|\alpha_M^\infty(A)\|.                              \tag{2.16}
\]
Thus \(\Omega_0\) is bounded with norm one.

\LS{1}{7}{Extend and identify the coordinate limits.}

By (1.3), \(\A_{\rm alg}\) is norm dense in \(\A_\infty\).  Therefore
\(\Omega_0\) extends uniquely and continuously to a state
\(\Omega_\infty\) on \(\A_\infty\).  Equations (2.6) and (2.10) give
(1.4).  \QedStep

\section{Spatial-cutoff interface}

\begin{corollary}[iterated cutoff then resolution limit]\label{cor:cutoff}
Suppose that for each \(N\) there is a directed cutoff family of states
\(\{\omega_{N,g}\}_g\subset S(\A_N)\) satisfying
\[
 \lim_{g\uparrow1}\omega_{N,g}(B)=\omega_N(B),\qquad B\in\A_N. \tag{3.1}
\]
Then the subnet obtained in Theorem \ref{thm:compactness} satisfies
\[
 \lim_i\left(\lim_{g\uparrow1}
 \omega_{N(i),g}(\alpha_M^{N(i)}(A))\right)
 =\Omega_\infty(\alpha_M^\infty(A))                            \tag{3.2}
\]
for every fixed \(M\) and \(A\in\A_M\).
\end{corollary}

\begin{proof}
For fixed \(i\), (3.1), applied to
\(B=\alpha_M^{N(i)}(A)\), makes the inner limit equal to
\(\omega_{N(i)}(\alpha_M^{N(i)}(A))\).  The outer limit is (1.4).
\end{proof}

\begin{remark}[no joint-limit inference]
Equation (3.2) is an iterated limit.  It does not imply a joint limit.  The scalar
array
\[
 a_{N,R}:=\mathbf1_{\{R<N\}}                                   \tag{3.3}
\]
satisfies \(\lim_{R\to\infty}a_{N,R}=0\) for every fixed \(N\), but the cofinal
diagonal \(R=N-1\) has \(a_{N,N-1}=1\).
\end{remark}

\section{Preserved compatible symmetries}

\begin{proposition}
Suppose \(\vartheta_N\in\operatorname{Aut}(\A_N)\) obey
\[
 \vartheta_N\alpha_M^N=\alpha_M^N\vartheta_M                  \tag{4.1}
\]
and \(\omega_N\vartheta_N=\omega_N\) for every \(N\).  Then there is an induced
\(\vartheta_\infty\in\operatorname{Aut}(\A_\infty)\), and the state from
Theorem \ref{thm:compactness} satisfies
\[
 \Omega_\infty\vartheta_\infty=\Omega_\infty.                  \tag{4.2}
\]
\end{proposition}

\begin{proof}
Fix \(M,A\).  Eventually \(N(i)\ge M\), and
\[
 \omega_{N(i)}(\alpha_M^{N(i)}(\vartheta_MA))
 =\omega_{N(i)}(\vartheta_{N(i)}(\alpha_M^{N(i)}A))
 =\omega_{N(i)}(\alpha_M^{N(i)}A).                             \tag{4.3}
\]
Taking the limit gives \(\rho_M\vartheta_M=\rho_M\).  Equation (4.2) first
follows on \(\A_{\rm alg}\) and then on \(\A_\infty\) by norm continuity.
\end{proof}

\section{Sharp limitations}

\begin{proposition}[the compactness conclusion cannot be upgraded formally]
The hypotheses of Theorem \ref{thm:compactness} imply none of the following:
\begin{enumerate}
\item convergence of the full sequence;
\item convergence in the norm of state functionals;
\item normality of the limiting state in a prescribed representation;
\item convergence of dynamics or the ground-state spectral condition.
\end{enumerate}
\end{proposition}

\begin{proof}
For (1), take \(\A_N=\C^2\), all maps equal to the identity, and alternate the two
point-evaluation states.

For (2)--(3), take the constant system \(\A_N=B(\ell^2(\mathbb N))\) and
\(\omega_n(A)=\langle e_n,Ae_n\rangle\).  Every weak-star cluster state \(\rho\)
annihilates compact operators because
\(\langle e_n,Ke_n\rangle\to0\) for every compact \(K\).  Hence \(\rho\) is
singular.  For \(P_n=|e_n\rangle\langle e_n|\) and
\(U_n=2P_n-I\),
\[
 \|U_n\|=1,\qquad \omega_n(U_n)=1,\qquad \rho(U_n)=-1.          \tag{5.1}
\]
Therefore \(\|\omega_n-\rho\|=2\), so no norm convergence occurs.

For (4), the hypotheses mention neither automorphism groups nor Hamiltonians.
Consequently no statement involving them is a logical consequence.  More concretely,
one may put arbitrary, mutually incompatible dynamics on the constant algebra without
changing any hypothesis or any step of the proof.
\end{proof}

\begin{remark}[subnet versus subsequence]
No separability assumption was made.  Thus weak-star compactness yields a subnet, not
necessarily a subsequence.  If every \(\A_M\) is separable, the state spaces are
weak-star metrizable on bounded sets, and a standard diagonal extraction gives a
subsequence.
\end{remark}

\end{document}
